\begin{abstract}
\noindent Surveillance comparisons of pathogen hospitalization rely on unexamined denominator conventions for unrecorded outcomes and overlook clustering and multiplicity. We estimated standardized hospitalization proportions among \textit{reported outbreak cases} (rather than intrinsic pathogen virulence) across 4,133 confirmed foodborne outbreaks in the US National Outbreak Reporting System (2009--2023), restricted to 12 shared setting--period cells (3,176 outbreaks; 72,132 cases; 90.4\% of weight in 2009--2018). Without assumptions on unrecorded hospitalization outcomes, 4 of 15 sample pairs separated (Shiga toxin--producing \textit{Escherichia coli} [STEC] above \textit{Campylobacter}, \textit{Clostridium perfringens}, and norovirus GI and GII). Nominal outbreak-clustered ($p = 0.024$) and exposure-state-clustered ($p = 0.033$) tests placed STEC above norovirus GII, whereas primary studentized bootstrap-$t$ inference rejected zero of 30 one-sided hypotheses ($p = 0.249$; all 95\% ranks $[1,6]$); thus worst-case statistical support is fragile and method-dependent. Within-cell hospitalization was higher in partially than fully recorded outbreaks ($+0.05$ to $+9.02$ percentage points [pp]), though STEC showed negative standardized observed-versus-full ($-2.53$ pp) and multistate ($-1.28$ pp) contrasts. Under bounded missingness departures anchored to external surveillance benchmarks, simultaneous 95\% ranks separated STEC $[1,1]$, \textit{Salmonella} $[2,2]$ and \textit{Campylobacter} $[3,3]$ above $[4,6]$ up to a joint studentized bootstrap-$t$ breakdown deviation of 11.6 pp (13.5 pp for \textit{Campylobacter}; 7.9 pp under the worst-case critical value), whereas lower-tier orderings failed below 5 pp.

\vspace{0.5em}\noindent\textbf{Keywords:} foodborne disease outbreaks; hospitalization; partial identification; breakdown frontiers; rank confidence sets
\end{abstract}
